% A l'original la figura és a la dreta dels càlculs de l'apartat a, amb «(0,2 p)» en vermell
% dins del requadre; s'ha esborrat del retall i s'escriu com a \punts abans de la figura.
\begin{solapartat}
Els components verticals s'anul·len per simetria

\medskip
\punts{0,2} $E_{1x} = E_{4x} = 8,99\times 10^{9}\cdot\dfrac{10^{-6}}{8}\,\dfrac{\sqrt{2}}{2} = 794,6\,N/C$

\smallskip
\punts{0,2} $E_{2x} = E_{3x} = 8,99\times 10^{9}\cdot\dfrac{2\times 10^{-6}}{8}\,\dfrac{\sqrt{2}}{2} = 1589,2\,N/C$

\smallskip
\punts{0,2} $E_C = 2\cdot 794,6 - 2\cdot 1589,2 = -1589,2\,N/C$

\smallskip
% A l'original falta el signe «=» en aquesta expressió.
\punts{0,2} $\vec{E}_C\ -1589,2\,\vec{\imath}\ \ N/C$

\medskip
\punts{0,2}
\figura[0.5]{sol_fig1.png}
\end{solapartat}

\begin{solapartat}
\punts{0,4} $V_C = 0$

\smallskip
\phantom{\punts{0,4}} $V_A = V_1 + V_2 + V_3 + V_4$

\medskip
\punts{0,4} $V_A = 8,99\times 10^{9}\left(\dfrac{10^{-6}}{2} - \dfrac{2\times 10^{-6}}{2} + \dfrac{2\times 10^{-6}}{\sqrt{20}} - \dfrac{10^{-6}}{\sqrt{20}}\right) = -2,48\times 10^{3}\,V$

\smallskip
\punts{0,2} $V_A - V_C = -2,48\times 10^{3}\,V$
\end{solapartat}
